\chapter{AIを使った文献調査}
\label{chap3}

\begin{chapterbox}
\textbf{【この章で学ぶこと】}
\begin{itemize}
\item AIツールを活用した効率的な文献調査ワークフローを習得する --- 4ステップの体系的な調査手法を実践できるようになる
\item 情報の正確性を批判的に検証する方法を身につける --- AIが生成する情報の限界を理解し、ファクトチェックを行える
\item 研究の新規性を言語化する技法を理解する --- 先行研究との差分を明確に説明できるようになる
\end{itemize}
\end{chapterbox}
\vspace{5mm}

%=========================================
\section{研究調査の課題}
\label{sec:3.1}
%=========================================

\subsection{情報爆発の時代}

現代の研究者が直面する最大の課題は、学術情報の爆発的な増加である。プレプリントサーバーarXiv\index{arXiv}の投稿数は、2010年には年間約7万件であったが、2023年には年間約19万件に達し、わずか13年で約2.7倍に増加した。PubMedに登録される生物医学系論文も、年間約100万件に達している。

この情報爆発は、特に大学院生にとって深刻な問題である。自分の研究テーマに関連する論文をすべて読むことは物理的に不可能であり、重要な先行研究を見落とすリスクが常につきまとう。

\subsection{キーワードの壁}

文献調査\index{ぶんけんちょうさ@文献調査}で見落としが発生する大きな原因の一つは、「キーワードの壁」である。同じ概念でも分野によって異なる用語が使われることがある。例えば、以下のような例がある。

\begin{itemize}
\item 「転移学習（Transfer Learning）」と「ドメイン適応（Domain Adaptation）」
\item 「特徴量エンジニアリング（Feature Engineering）」と「表現学習（Representation Learning）」
\item 「異常検知（Anomaly Detection）」と「外れ値検出（Outlier Detection）」と「新奇性検出（Novelty Detection）」
\end{itemize}

自分が知らない用語で記述されている重要な研究を見逃す可能性は、特に分野横断的な研究において顕著である。従来の検索エンジンはキーワード一致に依存するため、この問題を解決できなかった。

\subsection{時間的制約}

大学院生の日常は、授業、研究、TA業務、学会準備など多忙を極める。その中で英語論文を精読する時間を確保することは容易ではない。1本の論文を深く理解するのに数時間かかることも珍しくなく、関連論文を網羅的に調査するには膨大な時間が必要である。

特に英語が母語でない日本語話者にとって、英語論文の読解には追加の負荷がかかる。専門用語の理解、複雑な文構造の把握、文脈に依存した表現の解釈など、言語的な障壁が調査効率を低下させる。

\subsection{思考の停滞}

文献調査を進める中で、多くの大学院生が経験するのが「思考の停滞」である。先行研究を読めば読むほど、「すでに研究されている」「自分の研究に新規性がない」と感じてしまうことがある。

この問題の根本は、自分の研究の新規性\index{しんきせい@新規性}をうまく言語化できないことにある。研究の新規性は、先行研究との差分を明確にすることで初めて可視化される。しかし、その差分を的確に表現する技術は、経験を積まなければ身につかないものであった。AIはこの言語化のプロセスを支援する。

%=========================================
\section{文献調査AIツールの使い分け}
\label{sec:3.2}
%=========================================

\subsection{対話型AI --- ブレスト・壁打ちの相棒}

\textbf{主なツール:} Gemini Pro、ChatGPT（GPT-4o）、Claude

対話型AIは、研究のアイデア出しや方向性の検討に最適である。人間の指導教員や研究仲間と壁打ちするように、AIと対話することで思考を整理できる。

\textbf{長所:}
\begin{itemize}
\item 自然言語で曖昧な質問ができる
\item 対話を通じて段階的に思考を深められる
\item 異なる視点からの提案を得られる
\item 複雑な概念の説明を求められる
\end{itemize}

\textbf{短所:}
\begin{itemize}
\item 出典が明示されないことが多い
\item 存在しない論文を生成する場合がある（ハルシネーション）
\item 学習データの時期によって最新情報が反映されない
\item 回答の正確性を自分で検証する必要がある
\end{itemize}

\subsection{検索特化型AI --- 出典付き回答でサーベイの初動}

\textbf{主なツール:} Perplexity\index{Perplexity}

検索特化型AIは、ウェブ検索とAIの回答生成を組み合わせたツールである。回答に出典リンクが付与されるため、情報の検証が容易である。

\textbf{長所:}
\begin{itemize}
\item 回答に出典URLが付く
\item リアルタイムのウェブ情報を反映できる
\item サーベイの初期段階で全体像を掴むのに適している
\item 英語と日本語の両方で検索可能
\end{itemize}

\textbf{短所:}
\begin{itemize}
\item 深い分析や批判的検討は苦手
\item 出典の質にばらつきがある（学術論文以外の情報も混在）
\item 対話の文脈保持能力は対話型AIに劣る
\end{itemize}

\subsection{論文分析AI --- 深掘りのための専門ツール}

\textbf{主なツール:} NotebookLM\index{NotebookLM}、Elicit\index{Elicit}、Connected Papers\index{Connected Papers}、Semantic Scholar\index{Semantic Scholar}

論文分析に特化したAIツールは、学術文献の深い理解を支援する。

\begin{table}[htbp]
\centering
\caption{論文分析AIツールの比較}
\label{tab:paper_analysis_tools}
\begin{tabular}{lll}
\toprule
ツール名 & 特徴 & 最適な用途 \\
\midrule
NotebookLM & アップロードした文書をもとに質問応答 & 特定論文の深い理解 \\
Elicit & 学術論文の検索・要約に特化 & 体系的な文献レビュー \\
Connected Papers & 引用関係の可視化 & 研究分野のマッピング \\
Semantic Scholar & AI支援の論文検索エンジン & 関連論文の網羅的発見 \\
\bottomrule
\end{tabular}
\end{table}

\subsection{目的別フローチャート}

文献調査の目的に応じて、最適なツールの組み合わせは異なる。以下に目的別の推奨フローを示す。

\begin{itembox}[l]{目的別の推奨フロー}
\textbf{パターンA: 新しい研究テーマの探索}
\begin{enumerate}
\item 対話型AI（Claude/ChatGPT）でブレインストーミング
\item Perplexityで関連分野の概要を把握
\item Connected Papersで研究分野のマッピング
\end{enumerate}

\vspace{3mm}
\textbf{パターンB: 既存テーマの先行研究調査}
\begin{enumerate}
\item Perplexityで主要な先行研究を特定
\item Semantic ScholarまたはElicitで網羅的な検索
\item NotebookLMで重要論文を深掘り
\item 対話型AIで新規性の言語化
\end{enumerate}

\vspace{3mm}
\textbf{パターンC: 特定論文の理解}
\begin{enumerate}
\item NotebookLMに論文をアップロードして質問
\item 対話型AIに要約・解説を依頼
\item Connected Papersで関連研究を確認
\end{enumerate}
\end{itembox}

%=========================================
\section{4ステップの文献調査ワークフロー}
\label{sec:3.3}
%=========================================

\subsection{Step 1: キーワード拡張 --- AIで関連キーワードを発見する}

文献調査の第一歩は、検索に使うキーワードの拡張\index{きーわーどかくちょう@キーワード拡張}である。自分が知っているキーワードだけでは、関連する研究を見落とす可能性がある。AIを使って関連キーワードを網羅的に発見しよう。

\begin{promptbox}
\textbf{【基本プロンプト例：キーワード拡張】}

私は「[研究テーマ]」について文献調査を行いたいと考えています。この研究テーマに関連する検索キーワードを、以下の観点から網羅的にリストアップしてください。

\begin{enumerate}
\item 同義語・類義語（英語・日本語両方）
\item 上位概念・下位概念
\item 関連する手法名・アルゴリズム名
\item 関連する応用分野
\item この分野で頻出するキーワード
\end{enumerate}

それぞれのキーワードについて、簡単な説明も付けてください。
\end{promptbox}

\begin{promptbox}
\textbf{【具体例 --- 「グラフニューラルネットワークによる薬物相互作用予測」の場合】}

私は「グラフニューラルネットワークによる薬物相互作用予測」について文献調査を行いたいと考えています。この研究テーマに関連する検索キーワードを、以下の観点から網羅的にリストアップしてください。

\begin{enumerate}
\item 同義語・類義語（英語・日本語両方）
\item 上位概念・下位概念
\item 関連する手法名・アルゴリズム名
\item 関連する応用分野
\item この分野で頻出するキーワード
\end{enumerate}

それぞれのキーワードについて、簡単な説明も付けてください。
\end{promptbox}

\textbf{出力の評価と活用法:}

AIが生成したキーワードリストは、そのまま使うのではなく、以下の手順で評価・活用する。

\begin{enumerate}
\item \textbf{知っているキーワードの確認} --- リストに自分が既知のキーワードが含まれているか確認する。含まれていない場合、プロンプトの改善が必要
\item \textbf{未知のキーワードの学習} --- 知らなかったキーワードについて、追加で調べる
\item \textbf{キーワードの優先順位付け} --- 自分の研究テーマとの関連度が高い順に並べ替える
\item \textbf{検索クエリの構築} --- 優先度の高いキーワードを組み合わせて、Google ScholarやSemantic Scholarの検索クエリを作成する
\end{enumerate}

\begin{screen}
\textbf{ヒント:} キーワード拡張は一度で完了するものではない。調査を進める中で新たなキーワードが見つかったら、このステップに戻って検索範囲を広げよう。反復的な調査が網羅性を高める。
\end{screen}

\subsection{Step 2: 全体像を掴むサーベイ --- Perplexity等で俯瞰する}

キーワードが揃ったら、次は研究分野の全体像を掴むサーベイ\index{さーべい@サーベイ}を行う。ここではPerplexityのような検索特化型AIが効果的である。

\paragraph{英語プロンプトの重要性}
学術情報の大部分は英語で公開されている。そのため、サーベイのプロンプトは英語で入力することを強く推奨する。日本語で入力すると、日本語のウェブページが優先的に検索され、最新の学術論文が見落とされる可能性がある。

\begin{promptbox}
\textbf{【英語でのサーベイプロンプト例】}

What are the current state-of-the-art approaches for drug-drug interaction prediction using graph neural networks? Please provide:

\begin{enumerate}
\item A brief overview of the field (2--3 paragraphs)
\item Key methods and their characteristics (with citations)
\item Major datasets used in this field
\item Open challenges and future directions
\item Key research groups and their contributions
\end{enumerate}
\end{promptbox}

\textbf{出典リンクの確認方法:}

Perplexityが提示する出典は、必ず以下の手順で確認する。

\begin{enumerate}
\item 出典リンクをクリックして、実際のページが存在するか確認する
\item 論文の場合、DOIまたはarXiv IDが正しいか確認する
\item Google Scholarで論文タイトルを検索し、被引用数や掲載ジャーナルを確認する
\item 情報が一次情報源に基づいているか確認する（ブログ記事やニュース記事ではなく、論文原文を参照する）
\end{enumerate}

\begin{screen}
\textbf{ヒント:} Perplexityの「Academic」モードを使うと、学術論文を優先的に検索してくれる。サーベイの初期段階では、このモードを活用しよう。
\end{screen}

\subsection{Step 3: 注目論文の深掘り --- アブストラクトから構造的要約へ}

全体像を掴んだら、特に重要と思われる論文を深く読み込む。AIを活用して、効率的に論文の要点を抽出しよう。

\begin{promptbox}
\textbf{【4点要約テンプレート】}

以下の論文のアブストラクト（および可能であれば本文）を読み、次の4点について日本語で構造的に要約してください。

【論文タイトル】: [タイトルを入力]

【アブストラクト】: [アブストラクトを貼り付け]

\begin{enumerate}
\item \textbf{背景（Background）}: この研究が取り組む問題は何か？なぜ重要か？
\item \textbf{手法（Method）}: どのようなアプローチ・手法を提案しているか？技術的な新規性は何か？
\item \textbf{結果（Results）}: 主要な実験結果は何か？既存手法と比較してどの程度の改善が得られたか？
\item \textbf{限界（Limitations）}: この研究の限界や未解決の課題は何か？今後の研究の方向性として何が示唆されているか？
\end{enumerate}

また、この論文の新規性を一文で表現してください。
\end{promptbox}

\paragraph{複数論文の比較要約}
重要な論文が複数見つかった場合、AIを使って比較表を作成できる。

\begin{promptbox}
\textbf{【複数論文の比較要約プロンプト】}

以下の3つの論文について、比較表を作成してください。

論文1: [タイトル] --- [アブストラクト]

論文2: [タイトル] --- [アブストラクト]

論文3: [タイトル] --- [アブストラクト]

比較の軸:
\begin{itemize}
\item 解決しようとしている問題
\item 提案手法の特徴
\item 使用データセット
\item 主要な結果（定量的な指標を含む）
\item 限界と課題
\item 自分の研究テーマ「[テーマ]」との関連性
\end{itemize}
\end{promptbox}

\begin{screen}
\textbf{ヒント:} NotebookLMを使えば、論文PDFを直接アップロードして質問できる。アブストラクトだけでなく、実験設定の詳細や図表の解釈についても質問できるため、深い理解が得られる。
\end{screen}

\subsection{Step 4: 研究の新規性を言語化する}

文献調査の最終目標は、自分の研究の新規性を明確に言語化することである。このステップが最も重要であり、同時に最も難しいステップである。

\begin{promptbox}
\textbf{【新規性言語化のためのプロンプト】}

私の研究テーマは「[テーマ]」です。以下に、関連する先行研究とその限界をまとめました。

先行研究1: [タイトル] / 手法: [手法の概要] / 限界: [限界]

先行研究2: [タイトル] / 手法: [手法の概要] / 限界: [限界]

先行研究3: [タイトル] / 手法: [手法の概要] / 限界: [限界]

私が提案する手法は「[自分の手法の概要]」です。

以下の点について、整理を手伝ってください。
\begin{enumerate}
\item 先行研究に共通する未解決の課題は何か？
\item 私の提案手法は、それらの課題をどのように解決するか？
\item 私の研究の新規性を、学術論文のIntroductionに書くような形式で3--4文にまとめてください。
\item この新規性の主張に対して、査読者が指摘しそうな弱点は何か？
\end{enumerate}
\end{promptbox}

\paragraph{差分マトリクスの作成}

\begin{promptbox}
\textbf{【差分マトリクス作成プロンプト】}

以下の先行研究と私の研究について、差分マトリクス\index{さぶんまとりくす@差分マトリクス}を作成してください。

比較軸として以下を含めてください:
\begin{itemize}
\item 対象とする問題設定
\item 使用するデータの種類
\item モデルアーキテクチャ
\item 学習方法
\item 評価指標
\item 計算コスト
\item 実用性
\end{itemize}

各セルには、具体的な違いを簡潔に記述してください。自分の研究の優位点には$\bigstar$をつけてください。
\end{promptbox}

\begin{screen}
\textbf{ヒント:} 新規性の言語化は、AIの出力をそのまま使うのではなく、あくまで「たたき台」として活用すること。最終的な新規性の主張は、指導教員や研究仲間と議論して磨き上げることが重要である。AIはあくまで思考の補助であり、研究の方向性を決めるのは自分自身である。
\end{screen}

%=========================================
\section{情報の批判的検証}
\label{sec:3.4}
%=========================================

\subsection{AIが出力した論文情報のファクトチェック手法}

AIは非常に有用なツールであるが、その出力を無批判に信頼することは危険である。特に文献調査においては、以下のような問題が発生する可能性がある。

\begin{itembox}[l]{ファクトチェック\index{ふぁくとちぇっく@ファクトチェック}のチェックリスト}
\begin{itemize}
\item[$\square$] 論文のタイトルと著者名が正確か（Google Scholarで検索）
\item[$\square$] 掲載ジャーナル/学会名が正しいか
\item[$\square$] 発表年が正確か
\item[$\square$] 引用されている数値や結果が論文原文と一致するか
\item[$\square$] 論文の主張がAIの要約と整合しているか
\end{itemize}
\end{itembox}

\textbf{具体的なファクトチェック手順:}

\begin{enumerate}
\item \textbf{Google Scholarでの検索:} AIが挙げた論文のタイトルをそのままGoogle Scholarに入力し、実在する論文かどうかを確認する
\item \textbf{DOIの確認:} DOIが提示されている場合、doi.orgで解決できるか確認する
\item \textbf{著者のホームページ確認:} 主要著者の研究室ページや個人ページで、該当論文が掲載されているか確認する
\item \textbf{数値の照合:} 重要な実験結果の数値は、必ず論文原文で確認する
\end{enumerate}

\subsection{存在しない論文への対処（ハルシネーション問題）}

AIが存在しない論文を「作り上げる」ことは、ハルシネーション\index{はるしねーしょん@ハルシネーション}の典型例である。もっともらしいタイトル、著者名、ジャーナル名が生成されるが、実際には存在しない論文であるケースが少なくない。

\textbf{ハルシネーションを疑うべきサイン:}
\begin{itemize}
\item タイトルが非常に一般的で、どの分野にも当てはまりそうな表現
\item 有名な研究者の名前が著者に含まれているが、その研究者の専門分野と一致しない
\item 掲載ジャーナルが非常に有名だが、分野が異なる
\item 発表年が最新すぎる、または古すぎる
\item DOIやURLが提示されない
\end{itemize}

\begin{promptbox}
\textbf{【ハルシネーション対処のプロンプト例】}

先ほど挙げていただいた論文について確認したいのですが、以下の論文はGoogle Scholarで見つかりませんでした。

「[論文タイトル]」

この論文は実在しますか？もし存在しない場合、同様の内容を扱う実在する論文を教えてください。その際、Google Scholarで確認できるように、正確なタイトルと著者名を記載してください。
\end{promptbox}

\begin{screen}
\textbf{ヒント:} Perplexityは出典リンク付きで回答を生成するため、ハルシネーションのリスクが比較的低い。論文の実在性が重要な場合は、Perplexityで確認するか、直接Google ScholarやSemantic Scholarで検索しよう。
\end{screen}

\subsection{バイアスの認識}

AIの回答には、学習データに起因するバイアス\index{ばいあす@バイアス}が含まれる可能性がある。文献調査において注意すべきバイアスには以下がある。

\paragraph{英語圏バイアス}
AIは英語のデータで主に学習されているため、英語圏の研究が過度に強調される傾向がある。日本語や他の言語で発表された重要な研究が見落とされる可能性がある。

\paragraph{引用数バイアス}
被引用数が多い「有名な」論文が優先的に紹介される傾向があり、最新の研究や小規模な研究グループの重要な成果が見落とされる可能性がある。

\paragraph{分野バイアス}
AIの学習データに含まれる論文の分野分布に偏りがあるため、特定の分野（例: 機械学習、自然言語処理）の研究が過度に強調される可能性がある。

\paragraph{出版バイアス}
肯定的な結果を報告する論文が優先される傾向があり、否定的な結果や再現性の問題を報告する論文が見落とされる可能性がある。

これらのバイアスを認識した上で、AIの出力を批判的に評価することが重要である。

%=========================================
\section{文献管理とまとめ}
\label{sec:3.5}
%=========================================

\subsection{調査結果の整理法}

文献調査で得られた情報は、体系的に整理することで初めて研究に活用できる。AIを使って、調査結果を表形式で整理しよう。

\begin{promptbox}
\textbf{【文献比較表の作成プロンプト】}

以下の論文情報をもとに、比較表を作成してください。

[論文情報をリストアップ]

表の列には以下を含めてください:
\begin{itemize}
\item 論文（著者, 年）
\item 研究目的
\item 提案手法
\item データセット
\item 主要結果
\item 限界・課題
\item 自分の研究との関連
\end{itemize}
\end{promptbox}

整理のポイントとして、以下の項目を意識すること。

\begin{itemize}
\item \textbf{時系列での整理:} 研究の発展の流れを把握するため、発表年順に並べる
\item \textbf{手法別の整理:} 異なるアプローチを分類し、各手法の特徴を比較する
\item \textbf{問題設定別の整理:} 同じ問題に取り組む研究をグループ化する
\end{itemize}

\subsection{AIで文献レビューの草稿を作成する}

文献調査の結果を論文の「関連研究」セクションにまとめる作業も、AIの支援を受けることができる。

\begin{promptbox}
\textbf{【文献レビュー草稿作成プロンプト】}

以下の先行研究の情報をもとに、論文の``Related Work''セクションの日本語草稿を作成してください。

[先行研究の要約リスト]

以下の要件を満たしてください:
\begin{enumerate}
\item 研究を意味のあるカテゴリに分類して整理する
\item 各カテゴリ内で、研究の発展を時系列で説明する
\item 各研究の貢献と限界を簡潔に述べる
\item 最後に、これらの先行研究と本研究の関係を明確にする
\item 学術論文にふさわしい客観的な文体を使用する
\end{enumerate}
\end{promptbox}

\begin{screen}
\textbf{ヒント:} AIが生成した文献レビューの草稿は、あくまで出発点である。必ず自分で読み直し、以下を確認すること。(1) 論文の内容が正確に要約されているか、(2) 研究間の関係が正しく記述されているか、(3) 自分の研究の位置づけが明確か。
\end{screen}

\subsection{文献管理ツールとの連携}

AIで発見した論文は、文献管理ツール\index{ぶんけんかんりつーる@文献管理ツール}に登録して一元管理しよう。主要な文献管理ツールとその特徴は表\ref{tab:reference_tools}の通りである。

\begin{table}[htbp]
\centering
\caption{主要な文献管理ツールの比較}
\label{tab:reference_tools}
\begin{tabular}{lll}
\toprule
ツール名 & 特徴 & 料金 \\
\midrule
Zotero\index{Zotero} & オープンソース、ブラウザ拡張が便利 & 無料（300MBまで） \\
Mendeley\index{Mendeley} & PDF管理が優秀、Elsevierとの連携 & 無料（2GBまで） \\
Paperpile & Google Docsとの連携が強力 & 有料（学生割引あり） \\
EndNote\index{EndNote} & 多くの大学で機関ライセンスあり & 有料 \\
\bottomrule
\end{tabular}
\end{table}

\begin{itembox}[l]{AIと文献管理ツールの連携ワークフロー}
\begin{enumerate}
\item AIで文献調査を行い、重要な論文を特定する
\item Google ScholarやSemantic Scholarで論文を確認し、文献管理ツールに登録する
\item 文献管理ツールでタグ付け・メモの追加を行う
\item 論文執筆時に、文献管理ツールから引用を挿入する
\end{enumerate}
\end{itembox}

\begin{screen}
\textbf{ヒント:} ZoteroのブラウザExtensionを使えば、Google ScholarやarXivのページからワンクリックで論文を登録できる。AIで見つけた論文を効率的に管理するためにも、文献管理ツールの導入を強く推奨する。
\end{screen}

%=========================================
\section*{演習問題}
\addcontentsline{toc}{section}{演習問題}
%=========================================

\begin{enumerate}
\item \textbf{演習1: キーワード拡張の実践}

自分の研究テーマ（または興味のある研究トピック）について、AIを使ってキーワード拡張を行いなさい。以下の手順に従うこと。

\begin{enumerate}
\item 対話型AIに研究テーマを説明し、関連キーワードを少なくとも20個生成させる
\item 生成されたキーワードを「知っていたもの」と「知らなかったもの」に分類する
\item 知らなかったキーワードのうち、最も重要だと思うもの3つについて、Google Scholarで検索して関連論文を見つける
\item この演習を通じて発見した新しい視点や気づきをまとめる
\end{enumerate}

\item \textbf{演習2: ファクトチェックの実践}

以下の手順で、AIのハルシネーション検出を体験しなさい。

\begin{enumerate}
\item 対話型AIに「[自分の研究分野]における重要な論文を10本挙げてください」と質問する
\item 挙げられた10本の論文すべてについて、Google Scholarで実在を確認する
\item 存在しない論文があった場合、その数と割合を記録する
\item ハルシネーションが発生した論文について、どのような特徴があったかを分析する
\item 結果をもとに、AIの出力を検証する際の注意点をまとめる
\end{enumerate}

\item \textbf{演習3: 4点要約の実践}

自分の研究テーマに関連する英語論文を1本選び、以下を行いなさい。

\begin{enumerate}
\item まず自分でアブストラクトを読み、4点（背景・手法・結果・限界）を日本語でまとめる
\item 次に、本章のテンプレートを使ってAIに4点要約を依頼する
\item 自分の要約とAIの要約を比較し、(a) AIが正確に要約できた点、(b) AIが見落とした点、(c) AIが誤って解釈した点、をそれぞれ挙げる
\item この比較を通じて、AI要約を活用する際のベストプラクティスをまとめる
\end{enumerate}

\item \textbf{演習4: 新規性の言語化}

以下の手順で、自分の研究の新規性を言語化しなさい。

\begin{enumerate}
\item 自分の研究テーマに最も近い先行研究を3本特定する
\item 各先行研究について、手法と限界を整理する
\item 本章のプロンプトテンプレートを使って、AIに新規性の言語化を依頼する
\item AIの出力を批判的に検討し、修正・改善する
\item 最終的な新規性の記述を、指導教員やゼミのメンバーに共有してフィードバックを得る
\end{enumerate}

\item \textbf{演習5: 文献レビューの草稿作成}

自分の研究テーマについて、以下の手順で文献レビューの草稿を作成しなさい。

\begin{enumerate}
\item これまでの演習で収集した文献情報をもとに、AIに文献レビューの草稿を作成させる
\item 草稿を読み、以下の観点で評価する: (a) 分類は適切か、(b) 時系列の流れは正確か、(c) 各研究の記述は正確か
\item 問題があった箇所を修正し、自分の研究の位置づけを加筆する
\item 完成した草稿を文献管理ツール（Zoteroなど）の引用情報と照合し、参考文献リストを作成する
\end{enumerate}
\end{enumerate}

\vspace{5mm}
\begin{screen}
\textbf{本章のまとめ:} AIは文献調査を効率化する強力なツールであるが、万能ではない。AIの出力を鵜呑みにせず、常にファクトチェックを行い、批判的思考を維持することが不可欠である。文献調査の目的は、先行研究を知ることだけでなく、自分の研究の新規性と価値を明確にすることにある。AIを「思考の補助輪」として活用し、自分自身の研究力を高めてほしい。
\end{screen}
